% Generated draft: requires mathematical and source review.
\documentclass[11pt]{article}
\usepackage[margin=25mm]{geometry}
\usepackage{fontspec}
\setmainfont{Latin Modern Roman}
\usepackage{amsmath,amssymb,mathtools,unicode-math}
\setmathfont{Latin Modern Math}
\usepackage{enumitem,needspace,xurl,hyperref}
\setlength{\parindent}{0pt}
\setlength{\parskip}{5pt}
\setlength{\emergencystretch}{3em}
\newcommand{\N}{\mathbb{N}}
\newcommand{\Z}{\mathbb{Z}}
\newcommand{\Q}{\mathbb{Q}}
\newcommand{\R}{\mathbb{R}}
\newcommand{\C}{\mathbb{C}}
\newcommand{\F}{\mathbb{F}}
\newcommand{\A}{\mathbb{A}}
\newcommand{\PP}{\mathbb{P}}
\newcommand{\E}{\mathbb{E}}
\newcommand{\eps}{\varepsilon}
\newcommand{\dd}{\,\mathrm{d}}
\begin{document}
\section*{UnsolvedMath ID: 30004349}
\textbf{Source problem code:} catalogue entry
\subsection{Definitions and mathematical statement}
For a prime number \(p\), let
\[
\Z_{(p)}=\{a/b\in\Q:a,b\in\Z,\ b\ne0,\ p\nmid b\}.
\]
In the usual \(\mathrm{SL}_2(\Z)\) setting, a congruence subgroup is a subgroup \(\Gamma\leq\mathrm{SL}_2(\Z)\) containing
\[
\Gamma(N)=\ker\bigl(\mathrm{SL}_2(\Z)\longrightarrow\mathrm{SL}_2(\Z/N\Z)\bigr)
\]
for some integer \(N\geq1\). The supplied record does not explicitly fix the ambient group, so this conventional interpretation requires confirmation.

Under the resolution-based convention, a normal Noetherian surface \(S\) admitting a resolution has rational singularities if there is a proper birational morphism \(f:\widetilde S\to S\), with \(\widetilde S\) regular, such that
\[
f_*\mathcal O_{\widetilde S}=\mathcal O_S,
\qquad
R^i f_*\mathcal O_{\widetilde S}=0\quad\text{for every integer }i>0.
\]
The supplied record does not state the precise convention or hypotheses intended for rational singularities of the integral models in question.

\textbf{Classification problem.} Let \(\mathcal G\) denote the intended class of congruence subgroups, and let \(X_\Gamma\) denote the integral modular-curve model intended by the source for \(\Gamma\in\mathcal G\). Determine exactly which \(\Gamma\in\mathcal G\) have the property that \(X_\Gamma\) has rational singularities. More generally, determine exactly which pairs
\[
(p,\Gamma),\qquad p\text{ prime},\quad\Gamma\in\mathcal G,
\]
have the property that \((X_\Gamma)_{\Z_{(p)}}\) has rational singularities. If \(X_\Gamma\) is a scheme over \(\operatorname{Spec}\Z\), the displayed localization means
\[
(X_\Gamma)_{\Z_{(p)}}
=X_\Gamma\times_{\operatorname{Spec}\Z}\operatorname{Spec}\Z_{(p)}.
\]

\textbf{Unresolved qualifications.} The record does not identify the construction or base of \(X_\Gamma\), specify whether it is compactified, distinguish a coarse scheme from a moduli stack, or state normality and other hypotheses needed for the singularity condition. These choices must be recovered from the cited source before the classification target is unambiguous; no particular integral model is selected here.
\subsection{Short English statement}
Classify the congruence subgroups whose intended integral modular-curve models have rational singularities, and classify all prime--subgroup pairs for which the corresponding localized models have rational singularities. The precise model and singularity convention remain unspecified. The supplied record labels the problem partially solved; this is a supplied status assessment, not an independently verified conclusion.
\subsection{Sources}
\begin{itemize}
\item Supplied local UnsolvedMath registry record, permanent ID \(30004349\), source problem code \texttt{catalogue entry}, titled \emph{Rational Singularities of Integral Modular Curve Models}; collection: Oberwolfach Reports Open Problems. The supplied proposed year is 2019; no individual proposer is identified.
\item Supplied report citation: \emph{Modular Forms (2020)}, Oberwolfach Reports, pp.~3529--3606, DOI \texttt{10.4171/owr/2019/57}. Source link: \url{https://doi.org/10.4171/owr/2019/57}. The background describes Oberwolfach Reports as documenting mathematical workshops and their open-problem discussions.
\item Supplied literature reference: ``OWR report, Question 3 and known cases,'' \url{https://doi.org/10.4171/owr/2019/57}.
\item Literature assessment checked at \texttt{2026-08-22}: supplied dated curation, not independently verified. It labels the problem partially solved and reports positive answers in substantial cases, including the Raynaud range and additional cases from Manin-conjecture results, while reporting no general classification for all \(\Gamma\) and \(p\). The supplied background explicitly calls for independent verification of time-sensitive status and cited claims; it gives no precise hypotheses for these known cases.
\item The registry labels the statement \texttt{corrected\_verified}, explaining that only formula delimiters, notation, grammar, and typography were repaired, with mathematical content and scope preserved. This supplied label does not establish mathematical correctness or resolve the missing model specifications.
\end{itemize}
\end{document}
